\section{Experimental setup}\label{sec:setup}

\subsection{Scenes and cells}\label{sec:scenes}

The study uses \nScenes{} scenes (Table~\ref{tab:scenes}) from three acquisition groups: \emph{ground} scans, terrestrial and mobile laser scans of streets and heritage sites, which are three-dimensional and dense near the scanner; \emph{aerial} surveys of terrain and cities from aircraft and drones, which are 2.5-dimensional; and \emph{forest} scans, where vegetation fills the volume, together with a vegetated urban park. One photogrammetric object scan is used at one million points only. The scenes come from public benchmarks (Semantic3D, Paris-Lille-3D, Toronto-3D), national mapping agencies (PNOA in Spain, LINZ in New Zealand) and surveys of the authors' institution; Table~\ref{tab:scenes} gives the sources and the native point counts.

\begin{table*}[t]
\centering\footnotesize
\caption{The scenes. Flatness is the score of Section~\ref{sec:rule} (a range where it varies with size); cells counts the sizes at which a scene appears, from one million to one billion points. Every cell is a seeded strict subset of the scene's native points, so no cell contains interpolated points.}
\label{tab:scenes}
% generated by scripts/paper/tables.py; do not edit
\begin{tabular}{llllrrl}
\toprule
Scene & Group & Acquisition & Source & Native points & Flatness & Cells \\
\midrule
Alhambra & ground & TLS & own survey, Granada & 99\,999\,776 & 0.28 & 5 (1M--100M) \\
Cathedral & ground & TLS & Semantic3D & 95\,309\,849 & 0.45--0.51 & 4 (1M--50M) \\
DomFountain & ground & TLS & Semantic3D & 121\,466\,218 & 0.41--0.59 & 5 (1M--100M) \\
Lille & ground & MLS, Velodyne HDL-32E & Paris-Lille-3D & 120\,000\,000 & 0.48 & 4 (1M--100M) \\
Neugasse & ground & TLS & Semantic3D & 50\,122\,464 & 0.53 & 3 (1M--25M) \\
OttobrunnCampus & ground & TLS & -- & 111\,782\,240 & 0.56 & 4 (1M--100M) \\
ParisLux & ground & MLS, Velodyne HDL-32E & Paris-Lille-3D & 50\,000\,000 & 0.47 & 3 (1M--25M) \\
Toronto & ground & MLS, 32-line & Toronto-3D & 78\,320\,210 & 0.48 & 3 (1M--25M) \\
\midrule
Aneto & aerial & ALS & PNOA-LiDAR 2023 & 79\,717\,975 & 0.61 & 4 (1M--50M) \\
Barcelona & aerial & ALS & PNOA-LiDAR 2021 & 49\,267\,228 & 0.69 & 3 (1M--25M) \\
SanAndreas & aerial & UAV (senseFly eBee) & -- & 1\,191\,514\,616 & 0.53 & 7 (1M--1B) \\
SolarPlantation & aerial & UAV & own survey, solar plant & 728\,841\,014 & 0.81 & 4 (1M--100M) \\
Waikato Hamilton & aerial & ALS & LINZ Waikato LiDAR & 6\,015\,584\,388 & 0.64 & 7 (1M--1B) \\
\midrule
GaliciaForest & forest & TLS & 22 TLS plots, Galicia & 2\,052\,645\,430 & 0.50 & 8 (1M--1B) \\
GaliciaForest ES 05 & forest & TLS & one GaliciaForest plot & 114\,290\,378 & 0.52 & 5 (1M--100M) \\
Retiro & forest & ALS & PNOA-LiDAR 2024 & 39\,563\,361 & 0.89 & 3 (1M--25M) \\
\midrule
Museum & other & photogrammetry & Sketchfab & 2\,449\,789 & 0.25 & 1 (1M) \\
\bottomrule
\end{tabular}

\end{table*}

A \emph{cell} is one scene at one size. Sizes form a ladder from one million to one billion points, every tier a seeded strict subset of the one above, so that a scaling curve measures size alone and not a changing scene. The \nCells{} cells of the study (Table~\ref{tab:budgets}) are the \nSearchedCells{} that were searched, from one million to 250 million points, and \nTransferCells{} transfer cells at 500 million and one billion points, which reuse the schema selected on the same scene at 250 million and are measured like the rest. Cells up to 50 million points were measured on the workstation and cells of 100 million points and more on the server (Section~\ref{sec:hardware}); no ratio in this paper crosses the two machines.

\begin{table}[t]
\centering\footnotesize
\caption{Cells and budgets per size: the queries replayed from the held-out trace, and for the searched cells the genetic arm's generated schemas per generation, generations and population. The tuned-singles arm is the same 55 schemas at every size.}
\label{tab:budgets}
% generated by scripts/paper/tables.py; do not edit
\begin{tabular}{lrrll}
\toprule
Points & Cells & Queries & Arm & GA budget \\
\midrule
1M & 17 & 3000 & search & 24 $\times$ 2 $\times$ 16 \\
5M & 16 & 3000 & search & 24 $\times$ 2 $\times$ 16 \\
25M & 16 & 2000 & search & 24 $\times$ 2 $\times$ 12 \\
50M & 6 & 1500 & search & 32 $\times$ 3 $\times$ 20 \\
100M & 9 & 1000 & search & 32 $\times$ 3 $\times$ 20 \\
250M & 3 & 600 & transfer & -- \\
500M & 3 & 400 & transfer & -- \\
1B & 3 & 300 & transfer & -- \\
\bottomrule
\end{tabular}

\end{table}

\subsection{Query workloads}\label{sec:workloads}

The queries are the ones a point-cloud processing pipeline issues, not random points. For every cell, three stages of a standard pipeline were run over the cloud and their queries recorded: normal estimation, a $k$-nearest-neighbour query with $k = 16$ at every point; radius outlier removal, a radius query at every point with a radius of four times the cloud's mean point spacing; and Euclidean clustering, a radius query at 2.5 times the spacing. The mean spacing is computed exactly from the nearest-neighbour distances up to 130 million points and from point-anchored blocks above that. Each stage is subsampled uniformly to 20\,000 query centres, so a cell's trace holds 60\,000 queries, one third of them $k$-nearest-neighbour and two thirds radius, in the stage-major order the pipeline produced them.

Each trace is split in half, stratified by query type and parameter with a fixed seed and keeping the original order within each half. The \emph{opt} half is what the search sees and the \emph{test} half is what every number in this paper is measured on. A measurement replays the test half with a uniform stride down to the cell's query budget (Table~\ref{tab:budgets}), so the replayed queries stay spread over the whole cloud and consecutive queries are close in space, as they are in the pipeline. Query centres are stored in world coordinates and transformed to the cloud's local frame in double precision at load time.

Section~\ref{sec:composition} also reports seven variants of this workload with larger radii, larger $k$, and axis-aligned range and count queries built on the same centres, and Section~\ref{sec:ooc} adds a workload of 30 large boxes per cloud for the out-of-core regimes.

\subsection{Measurement protocol}\label{sec:protocol}

Every process is pinned to one core and, on the server, to the memory of that core's socket. Building is timed separately and excluded from every query latency. A schema is warmed up by one untimed pass over the replayed queries before its timed passes, and the timed passes run with telemetry compiled out. A schema's latency is the mean per-query time of a pass; a cell's result for a schema is the median of that mean over three processes that run the schemas in rotated order, so that a slow stretch of the machine hits every schema in turn. Two schemas within 3\% of each other are reported as a tie.

Two timing protocols are used, and the paper says which one each table rests on. In the \emph{interleaved} protocol all indexes of a process are built first and each pass runs one sweep per schema, starting one schema later each pass; this is the protocol of the search and of the held-out e9 numbers, and it is fair between the schemas of one process because they share the same conditions. It is, however, slower than timing a schema alone: with several large indexes resident, a memory-bound tree runs up to about a fifth slower, and Section~\ref{sec:mechanism} shows that most of that penalty is the reach of the translation look-aside buffer, since it nearly disappears on two-megabyte pages. In the \emph{sequential} protocol a schema is built, warmed up and timed before the next is built, which is within a few percent of timing it alone. The comparison with external libraries (Section~\ref{sec:libraries}), whose drivers time one setting after another, and the headline panel at 100 million points (Section~\ref{sec:levers}) use the sequential protocol; ratios are never taken across protocols.

Every pass recomputes the per-query counts and an order-independent hash of the returned identifiers and compares them with the first pass; the same hash is computed by the external drivers, so a library's answers are checked against the engine's query by query. The checks are reported with the results.

\subsection{External libraries}\label{sec:drivers}

Five libraries are timed on the same clouds and traces through one driver protocol: nanoflann~\cite{blanco2014nanoflann} (a sliding-midpoint kd-tree, leaf 10 by default, 16 to 256 swept), the Point Cloud Library~\cite{rusu2011pcl} (its FLANN kd-tree~\cite{muja2009flann}, and its octree search at voxel sizes from the extent over 128 to over 4096), Open3D~\cite{zhou2018open3d} (a FLANN kd-tree over a double-precision matrix), the octree of Behley et al.~\cite{behley2015octree} (bucket 32 by default, 8 to 512 swept; radius queries only), and the R-tree of Boost.Geometry with sort-tile-recursive bulk loading (16 entries per node as documented, 8 to 64 swept). Each driver reads the same cloud in the same local single-precision frame, replays the same stride of the same trace, times each search call alone around the library's own function, warms up with one pass and times ten, and runs in three rounds with the process order rotated. The reported latency is the median over rounds.

\subsection{Hardware}\label{sec:hardware}

The server has two Intel Xeon Gold 6140 processors (18 cores and 24.75\,MB of last-level cache per socket) and 251\,GB of memory, runs Linux with transparent huge pages in \emph{madvise} mode, and compiles the engine with GCC~12; it also holds an NVIDIA Quadro RTX 6000 used for the GPU measurements of Section~\ref{sec:composition}. The workstation has an Intel Core i7-14700KF (a hybrid processor; runs are pinned to a performance core) and 64\,GB of memory, runs Windows~11 and compiles the engine with MSVC. Cells up to 50 million points ran on the workstation and cells of 100 million points and more, the external libraries, the out-of-core and the layout experiments on the server. The server was shared with other users' jobs on other cores throughout; the rotated processes and medians of Section~\ref{sec:protocol} are there to absorb that.
